% =============================================================
%  Chapter 7 — Persuasion Bounds and Propaganda
%  Lean source: FormalAnthropology/IDT_Signal_Ch07.lean
% =============================================================
\chapter{Persuasion Bounds and Propaganda}
\label{ch:persuasion}

\epigraph{%
  ``The ideas of the ruling class are in every epoch the ruling ideas:
  i.e., the class which is the ruling material force of society is at the
  same time its ruling intellectual force.''}%
  {Karl Marx, \textit{The German Ideology} (1846)}

% ----------------------------------------------------------------
\section{Introduction: The Algebra of ``Hearts and Minds''}
\label{sec:pers-intro}

Every propagandist, preacher, and politician confronts the same constraint,
whether they know it or not: \emph{you cannot add more complexity to someone's
worldview than the complexity of your message allows}.  A bumper sticker
cannot produce a Ph.D.-level understanding; a three-word slogan cannot
install a twenty-variable causal model.  Conversely, a complex treatise
cannot simplify a mind below its current depth---complexity is added, not
subtracted, by composition.

Antonio Gramsci grasped this intuitively when he argued that hegemony requires
not just force but \emph{cultural production}---the sustained manufacture of
complex ideological structures capable of reshaping the depth of popular
consciousness.\footnote{Gramsci, \textit{Prison Notebooks}, Q.~12 (1932).
The concept of ``cultural hegemony'' names exactly the problem that this
chapter formalizes: how much ideological depth can a ruling class actually
install in a population?}  Michel Foucault extended the insight: power
operates not by loud commands but by quiet restructuring of the categories
through which subjects perceive the world---what he called
\emph{power/knowledge}.\footnote{Foucault, \textit{Discipline and Punish}
(1975); \textit{The History of Sexuality}, vol.~1 (1976).}

This chapter provides the mathematical backbone for these intuitions.  We
define \emph{marginal depth}---how much a signal adds to a receiver's
interpretive complexity---and prove that it is bounded above by the signal's
own depth (the \emph{fundamental persuasion bound}).  We then show that
chaining messages is \emph{subadditive}: the total persuasive capacity of
a compound message is at most the sum of its parts.  Repetition helps, but
only linearly---never exponentially.  Even totalitarian propaganda, with
unlimited repetition, faces a linear ceiling.

These results have sharp implications for anthropological theory.  They
explain why ``hearts and minds'' campaigns fail not because of insufficient
repetition, but because of the fundamental \emph{subadditivity} of ideatic
composition.  You cannot shout your way to depth.


% ================================================================
\section{Core Definitions}
\label{sec:pers-defs}

\begin{definition}[Marginal Depth]
\label{def:marginal-depth}
The \emph{marginal depth} of signal $s$ for receiver $r$ is the increase
in depth caused by composition:
\[
  \mathrm{marginalDepth}(r,s) \;=\; \depth(r\compose s) - \depth(r),
\]
where $-$ denotes truncated (natural number) subtraction.
\end{definition}

The marginal depth of a sermon to a parishioner is how much \emph{more}
complex their worldview becomes after hearing it, relative to what they
had before.  It is inherently receiver-dependent: the same lecture may add
five ``layers'' to a novice and zero to an expert.

\begin{definition}[Persuasion Bound]
\label{def:persuasion-bound}
The \emph{persuasion bound} of a signal $s$ is its own depth:
\[
  \mathrm{persuasionBound}(s) \;=\; \depth(s).
\]
This is an \emph{absolute} ceiling, independent of the receiver.
\end{definition}

\begin{lstlisting}[caption={Marginal depth and persuasion bound in Lean~4.}]
def marginalDepth (r s : I) : ℕ :=
  IdeaticSpace.depth (IdeaticSpace.compose r s) -
    IdeaticSpace.depth r

def persuasionBound (s : I) : ℕ := IdeaticSpace.depth s
\end{lstlisting}


% ================================================================
\section{The Fundamental Persuasion Bound}
\label{sec:pers-fundamental}

\begin{theorem}[Marginal Depth Bound (7.1) --- The Fundamental Theorem of Propaganda]
\label{thm:marginal-depth-bound}
For all receivers $r$ and signals $s$:
\[
  \mathrm{marginalDepth}(r,s) \;\le\; \mathrm{persuasionBound}(s)
  \;=\; \depth(s).
\]
\end{theorem}

\begin{proof}[Proof sketch]
By subadditivity of composition depth: $\depth(r\compose s)\le \depth(r)+\depth(s)$.
Subtracting $\depth(r)$ from both sides gives $\mathrm{marginalDepth}(r,s)\le\depth(s)$.
\end{proof}

No matter how receptive the audience, a message cannot add more complexity
than it intrinsically contains.  A bumper sticker cannot produce a
Ph.D.-level understanding.  This is the \textbf{fundamental theorem of
propaganda theory}: persuasive capacity is bounded by signal depth.

The result vindicates Gramsci's insistence that effective hegemony requires
\emph{organic intellectuals}---agents capable of producing signals of
sufficient depth to restructure popular consciousness.  Simple slogans
(low depth) have low persuasion bounds.  Only sustained, complex cultural
production can deeply transform a worldview.

\begin{theorem}[Void Has Zero Marginal Depth (7.2)]
\label{thm:void-marginal-depth}
$\mathrm{marginalDepth}(r, \void) = 0$ for all $r$.
\end{theorem}

\begin{proof}[Proof sketch]
$r\compose\void = r$ by the right-identity axiom, so $\depth(r\compose\void)
- \depth(r) = 0$.
\end{proof}

Silence cannot persuade.  Quaker silence, Buddhist \emph{noble silence}, and
the minute of silence for the dead change nothing about the listener's
depth.  They are powerful \emph{socially}---as markers of solidarity, grief,
or contemplation---but informationally null.%
\footnote{This creates a sharp distinction between the \emph{social} and
\emph{informational} functions of silence, a distinction that Foucault's
analysis of ``incitement to discourse'' (\textit{History of Sexuality},
vol.~1) explores from the opposite direction.}

\begin{theorem}[Void Has Zero Persuasion Bound (7.3)]
\label{thm:void-persuasion-bound}
$\mathrm{persuasionBound}(\void) = 0.$
\end{theorem}

Independent of the receiver, the null message has zero persuasive capacity.
This is a property of the signal itself, not of any particular audience.

\begin{theorem}[Persuasion Bound Nonnegativity (7.4)]
\label{thm:persuasion-nonneg}
$\mathrm{persuasionBound}(s) \ge 0$ for all $s$.
\end{theorem}

Trivially true for natural numbers, but conceptually important: there are
no ``anti-signals'' that reduce the maximum possible effect below zero.
Every signal has non-negative persuasive potential.


% ================================================================
\section{Deeper Signals Are More Persuasive}
\label{sec:pers-monotonicity}

\begin{theorem}[Monotonicity of Persuasion Bound (7.5)]
\label{thm:persuasion-mono}
If $\depth(s_1) \le \depth(s_2)$, then
$\mathrm{persuasionBound}(s_1) \le \mathrm{persuasionBound}(s_2)$.
\end{theorem}

A complex ritual has \emph{more potential} to transform a participant's
worldview than a simple one.  This doesn't mean it \emph{will}---marginal
depth depends on the receiver---but it \emph{can}.  Complexity is a
\emph{necessary} (not sufficient) condition for deep persuasion.

This result gives formal content to Foucault's concept of
\emph{power/knowledge}: the capacity to restructure categories of thought
(persuasion) scales with the complexity of the discursive apparatus
(signal depth).  Simple commands have low persuasion bounds; complex
disciplinary regimes have high ones.

\begin{theorem}[Self-Composition Persuasion Bound (7.6)]
\label{thm:self-compose-persuasion}
$\mathrm{persuasionBound}(s\compose s) \le 2\cdot\mathrm{persuasionBound}(s).$
\end{theorem}

\begin{proof}[Proof sketch]
$\depth(s\compose s)\le \depth(s)+\depth(s) = 2\depth(s)$ by subadditivity.
\end{proof}

Repeating a message doubles the \emph{maximum possible} effect, but the
\emph{actual} effect depends on the receiver.  Propaganda works by
repetition, but with diminishing returns.

\begin{theorem}[ComposeN Persuasion Bound (7.7)]
\label{thm:composeN-persuasion}
$\mathrm{persuasionBound}(\compN{s}{n}) \le n\cdot\mathrm{persuasionBound}(s).$
\end{theorem}

Even with unlimited repetition, persuasive capacity grows at most
linearly---never exponentially.  This is why totalitarian propaganda,
despite its volume, has bounded effectiveness.  The 20th century
demonstrated this repeatedly: neither Nazi Germany nor the Soviet Union
achieved total ideological control despite saturation-level messaging.%
\footnote{Hannah Arendt's \textit{Origins of Totalitarianism} (1951)
documents both the ambition and the failure of totalitarian propaganda
to achieve total depth penetration.}

\begin{lstlisting}[caption={ComposeN persuasion bound in Lean~4.}]
theorem composeN_persuasion_bound (s : I) (n : ℕ) :
    persuasionBound (composeN s n) <=
      n * persuasionBound s
\end{lstlisting}


% ================================================================
\section{Subadditivity: Chaining Doesn't Multiply}
\label{sec:pers-subadditivity}

\begin{theorem}[Persuasion Subadditivity (7.8)]
\label{thm:persuasion-subadditivity}
\[
  \mathrm{persuasionBound}(s_1\compose s_2)
  \;\le\;
  \mathrm{persuasionBound}(s_1) + \mathrm{persuasionBound}(s_2).
\]
\end{theorem}

\begin{proof}[Proof sketch]
Direct from the subadditivity of depth under composition:
$\depth(s_1\compose s_2)\le\depth(s_1)+\depth(s_2)$.
\end{proof}

A two-part message cannot be more persuasive than the sum of its parts.
This rules out ``synergistic propaganda''---the idea that combining
messages could produce super-additive persuasion.  It is the formal
foundation of bounded rationality in cultural transmission.

The subadditivity theorem has immediate policy implications.  If a state
seeks to increase persuasive capacity, it must increase the \emph{depth} of
its messages, not merely their \emph{number}.  Quantity is strictly
dominated by quality as a persuasion strategy.

\begin{theorem}[Three-Message Persuasion Bound (7.9)]
\label{thm:three-message-bound}
\[
  \mathrm{persuasionBound}\bigl((s_1\compose s_2)\compose s_3\bigr)
  \;\le\;
  \mathrm{persuasionBound}(s_1) +
  \mathrm{persuasionBound}(s_2) +
  \mathrm{persuasionBound}(s_3).
\]
\end{theorem}

Multi-act rituals (liturgy with readings, sermon, and communion) have bounded
total impact $\le$ the sum of individual act complexities.  The extension to
$n$~messages follows by induction.

\begin{theorem}[Marginal Depth of Void Receiver (7.10)]
\label{thm:marginal-void-receiver}
$\mathrm{marginalDepth}(\void, s) = \depth(s).$
\end{theorem}

\begin{proof}[Proof sketch]
$\void\compose s = s$ by left-identity, and $\depth(\void)=0$, so
$\depth(s) - 0 = \depth(s)$.
\end{proof}

The ``blank slate'' receiver (\emph{tabula rasa}) absorbs every signal
fully.  This is the formal version of Locke's epistemology---and its
anthropological critique: no real human is a blank slate, so no real human
absorbs a signal fully.  The void receiver is a theoretical limit, never
achieved in practice.


% ================================================================
\section{Persuasion Through Chains}
\label{sec:pers-chains}

What happens when a signal passes through an intermediary?  A teacher
receives a text, processes it through their own background, and transmits
the result to a student.  How much depth can this chain deliver?

\begin{theorem}[Chain Persuasion Bound (7.11)]
\label{thm:chain-persuasion}
$\mathrm{marginalDepth}(r, r_1\compose s) \le \depth(r_1) + \depth(s).$
\end{theorem}

A teacher cannot transmit more depth than their own understanding plus the
source material's complexity.  Great teaching requires both deep knowledge
\emph{and} rich source texts.  This formalizes what every graduate student
knows intuitively: a mediocre advisor with a brilliant text, or a brilliant
advisor with a mediocre text, both produce bounded results.

\begin{theorem}[Marginal Depth of Self-Composition (7.12)]
\label{thm:marginal-self}
$\mathrm{marginalDepth}(r, r) \le \depth(r).$
\end{theorem}

Reflecting on your own ideas (self-composition) cannot add more than your
existing complexity.  Introspection has bounded returns.  This result
anticipates Chapter~\ref{ch:echo}'s analysis of echo chambers: if
self-reflection is bounded, then \emph{iterated} self-reflection is
bounded too.

\begin{theorem}[Marginal Depth Triangle Inequality (7.13)]
\label{thm:marginal-triangle}
$\mathrm{marginalDepth}(r, s_1\compose s_2) \le \depth(s_1) + \depth(s_2).$
\end{theorem}

The marginal impact of a compound message is bounded by the sum of its
parts' complexities.  No synergy, no multiplication---just addition.  This
is the triangle inequality for marginal depth, and it gives the persuasion
space a quasi-metric structure.


% ================================================================
\section{Propaganda Effectiveness Bounds}
\label{sec:pers-propaganda}

We now turn to the specific problem of propaganda: what happens when the
same message is repeated many times?

\begin{theorem}[Bounded Propaganda Effectiveness (7.14)]
\label{thm:bounded-propaganda}
After composing receiver $r$ with $n$ repetitions of signal $s$:
\[
  \depth\bigl(r\compose\compN{s}{n}\bigr)
  \;\le\;
  \depth(r) + n\cdot\depth(s).
\]
\end{theorem}

\begin{proof}[Proof sketch]
Combine the subadditivity of composition depth with the bound
$\depth(\compN{s}{n})\le n\cdot\depth(s)$ from the composeN depth theorem.
\end{proof}

Propaganda saturation---hearing the same message $n$ times---increases
depth at most linearly.  No amount of repetition can produce exponential
radicalization.  This is why ``deprogramming'' is possible: the damage is
bounded, and bounded damage can be reversed.%
\footnote{The bounded-propaganda theorem provides a formal counterargument
to the popular ``brainwashing'' hypothesis.  See Lifton,
\textit{Thought Reform and the Psychology of Totalism} (1961), for the
empirical evidence that even extreme indoctrination has bounded effects.}

\begin{theorem}[Resonant Propaganda (7.15)]
\label{thm:resonant-propaganda}
If $s_1\res s_2$, then $r\compose s_1 \res r\compose s_2$ for all~$r$.
\end{theorem}

Competing narratives that ``resonate'' with the same cultural themes produce
resonant (structurally similar) effects on receivers.  When Fox News and
MSNBC draw on the same patriotic imagery, they produce resonant---not
identical---reactions in their audiences.  The content differs; the
\emph{structure} of the interpretive effect is preserved.

\begin{theorem}[Double Propagandist Bound (7.16)]
\label{thm:double-propagandist}
$\mathrm{marginalDepth}(r, s_1\compose s_2) \le
\mathrm{persuasionBound}(s_1) + \mathrm{persuasionBound}(s_2).$
\end{theorem}

Allied propaganda campaigns---joint military-religious messaging,
corporate-governmental public relations---cannot exceed the sum of their
individual complexities.  Coalition propaganda doesn't synergize; it merely
adds.

\begin{theorem}[Atomic Signal Minimal Persuasion (7.17)]
\label{thm:atomic-persuasion}
If $s$ is atomic, then $\mathrm{persuasionBound}(s) \le 1$.
\end{theorem}

Single words, simple gestures, basic symbols---all have persuasion bound at
most~1.  To deeply transform a worldview, you need \emph{structured}
complexity, not atomic simplicity.  A thumbs-up cannot substitute for a
treatise.%
\footnote{This result formalizes the linguistic insight that single lexemes
have bounded semantic impact.  The power of language lies in
\emph{compositional} complexity, not lexical richness.}

\begin{theorem}[Iterated Propaganda Marginal Decay (7.18)]
\label{thm:iterated-propaganda}
$\mathrm{marginalDepth}(r, \compN{s}{n}) \le n\cdot\mathrm{persuasionBound}(s).$
\end{theorem}

\begin{proof}[Proof sketch]
By the composition depth bound and the composeN depth bound:
$\depth(r\compose\compN{s}{n})\le\depth(r)+n\cdot\depth(s)$,
so $\mathrm{marginalDepth}\le n\cdot\depth(s)$.
\end{proof}

The $k$-th repetition of propaganda is no more effective than the first---the
bound is uniform.  Each repetition has the same ceiling on marginal depth,
but the \emph{actual} marginal depth can only decrease in realistic models.
This formalizes the ``diminishing returns'' of propaganda.

\begin{lstlisting}[caption={Iterated propaganda bound in Lean~4.}]
theorem iterated_propaganda_bound (r s : I) (n : ℕ) :
    marginalDepth r (composeN s n) <=
      n * persuasionBound s
\end{lstlisting}


% ================================================================
\section{Conclusion}
\label{sec:pers-conclusion}

This chapter established the fundamental limits of persuasion within IDT\@.
The central result---the marginal depth bound
(Theorem~\ref{thm:marginal-depth-bound})---states that no signal can add
more complexity than it contains.  The subadditivity theorem
(\ref{thm:persuasion-subadditivity}) rules out synergistic propaganda, and
the bounded propaganda theorem (\ref{thm:bounded-propaganda}) caps the
effect of repetition at linear growth.

Together, these results provide formal foundations for Gramsci's theory of
cultural hegemony (persuasion requires complex cultural production),
Foucault's power/knowledge (persuasive capacity scales with discursive
depth), and the empirical observation that totalitarian propaganda has
bounded effectiveness.

Chapter~\ref{ch:lossy} turns from the \emph{sender's} constraints to the
\emph{receiver's}: what happens when the reception process itself is lossy?
